\clearpage
\SourceLeaf{179}%
\pdfbookmark[0]{Göttingische gelehrte Anzeigen. 1831 April 23.}{notice1831april23}
\begin{center}
\rule{.82\linewidth}{.35pt}\par
Göttingische gelehrte Anzeigen.\quad 1831 April 23.\par
\rule{.82\linewidth}{.35pt}
\end{center}
Eine am 15. April von dem Hofr. \textsc{Gauss} der Königl. Societät überreichte Vorlesung:
\begin{center}\itshape
Theoria residuorum biquadraticorum, commentatio secunda,
\end{center}
ist die Fortsetzung der bereits im sechsten Bande der \emph{Commentationes novae} abgedruckten Abhandlung, wovon auch in unsern Blättern zu seiner Zeit 1825 April 11 eine Anzeige gemacht war. Auch diese Fortsetzung, obgleich mehr als doppelt stärker wie die erste Abhandlung, erschöpft den überaus reichhaltigen Gegenstand noch nicht, und erst einer künftigen dritten Abhandlung wird die Vollendung des Ganzen vorbehalten bleiben.

Obgleich die Grundbegriffe dieser Lehren und der Inhalt der ersten Abhandlung als allen, die aus der höhern Arithmetik ein Studium gemacht haben, bekannt vorausgesetzt werden können, wollen wir doch jene zur Bequemlichkeit solcher Freunde dieses Theils der Mathematik, welchen die erste Abhandlung nicht gleich zur Hand ist, hier kurz in Erinnerung bringen. In Beziehung auf eine beliebige ganze Zahl $p$ heisst eine andere $k$ ein biquadratischer Rest, wenn es Zahlen der Form $x^4-k$ gibt, die durch $p$ theilbar sind; im entgegengesetzten Fall heisst sie biquadratischer Nichtrest von $p$. Es ist zureichend, sich hiebei auf den Fall einzuschränken, wo $p$ eine Primzahl der Form $4n+1$, und $k$ durch
\SourceLeaf{180}%
dieselbe nicht theilbar ist, da alle andere Fälle entweder für sich klar, oder auf diesen zurückzuführen sind.

Für einen solchen \emph{gegebenen} Werth von $p$ zerfallen sämmtliche durch $p$ nicht theilbare Zahlen in \emph{vier} Classen, wovon die eine die biquadratischen Reste, eine zweite solche biquadratische Nichtreste, die quadratische Reste von $p$ sind, enthält, und in die beiden übrigen die biquadratischen Nichtreste, welche zugleich quadratische Nichtreste sind, vertheilt werden. Das Princip dieser Vertheilung besteht darin, dass allemal entweder $k^n-1$, oder $k^n+1$, oder $k^n-f$, oder $k^n+f$ durch $p$ theilbar sein wird, wo $f$ eine ganze Zahl bedeutet, die $ff+1$ durch $p$ theilbar macht. Jeder, dem die elementarische Terminologie bekannt ist, sieht von selbst, wie diese Worterklärungen in dieselbe eingekleidet werden.

Die Theorie dieser Classificirung nicht nur für den an der Oberfläche liegenden Fall $k=-1$, sondern auch für die, subtile Hülfsuntersuchungen erfordernden, Fälle $k=\pm2$, findet sich in der ersten Abhandlung ganz vollendet. Im Anfang der gegenwärtigen Abhandlung wird nun zu grösseren Werthen von $k$ fortgeschritten: man braucht aber dabei zunächst nur solche in Betracht zu ziehen, die selbst Primzahlen sind, und der Erfolg zeigt, dass die Resultate am einfachsten ausfallen, wenn man die Werthe positiv oder negativ nimmt, je nachdem sie, absolut betrachtet, von der Form $4m+1$ oder $4m+3$ sind. Die Induction gibt hier sofort mit grosser Leichtigkeit eine reiche Ernte von neuen Lehrsätzen, wovon wir hier nur ein paar anführen. Die Numerirung der Classen mit $1,2,3,4$ wird auf die Fälle bezogen, wo $k^n$ den Zahlen $1,f,-1,-f$ congruent wird; zugleich ist für die Zahl $f$ immer derjenige Werth angenommen, welcher $a+bf$ durch $p$ theilbar macht, wenn $aa+bb$ die Zerlegung von $p$ in ein ungerades und ein gerades Quadrat vorstellt. So findet sich durch die Induction, dass die Zahl $-3$ allemal zu der Classe $1,2,3,4$ gehört, je nachdem $b,a+b,a,a-b$ durch $3$ theilbar ist; dass die Zahl $+5$ der Reihe nach zu jenen Classen gehört, je nachdem $b,a-b,a,a+b$ durch $5$ theilbar ist; dass die Zahl $-7$ in die Classe $1$ fällt, wenn $a$ oder $b$; in die Classe $2$, wenn $a-2b$ oder $a-3b$; in die Classe $3$, wenn $a-b$ oder $a+b$; in die Classe $4$, wenn $a+2b$ oder $a+3b$ durch $7$ theilbar ist. Aehnliche Theoreme ergeben sich in Beziehung auf die Zahlen $-11,+13,+17,-19,-23$ u. s. f. So leicht sich aber alle dergleichen specielle Theoreme durch die Induction entdecken lassen, so schwer scheint
\SourceLeaf{181}%
es, auf diesem Wege ein allgemeines Gesetz für diese Formen aufzufinden, wenn auch manches Gemeinschaftliche bald in die Augen fällt, und noch viel schwerer ist es, für diese Lehrsätze die Beweise zu finden. Die für die Zahlen $+2$ und $-2$ in der ersten Abhandlung gebrauchten Methoden vertragen hier keine Anwendung mehr, und wenn gleich andere Methoden ebenfalls das, was sich auf die erste und dritte Classe bezieht, zu erledigen dienen könnten, so zeigen sich doch solche zur Begründung von \emph{vollständigen} Beweisen untauglich.

Man erkennt demnach bald, dass man in dieses reiche Gebiet der höhern Arithmetik nur auf ganz neuen Wegen eindringen kann. Der Verf. hatte schon in der ersten Abhandlung eine Andeutung gegeben, dass dazu eine eigenthümliche Erweiterung des ganzen Feldes der höhern Arithmetik wesentlich erforderlich ist, ohne damals sich näher darüber zu erklären, worin dieselbe bestehe: die gegenwärtige Abhandlung ist dazu bestimmt, diesen Gegenstand ins Licht zu setzen.

Es ist dieses nichts anders, als dass für die wahre Begründung der Theorie der biquadratischen Reste das Feld der höhern Arithmetik, welches man sonst nur auf die reellen ganzen Zahlen ausdehnte, auch über die imaginären erstreckt werden, und diesen das völlig gleiche Bürgerrecht mit jenen eingeräumt werden muss. Sobald man dies einmal eingesehen hat, erscheint jene Theorie in einem ganz neuen Lichte, und ihre Resultate gewinnen eine höchst überraschende Einfachheit.

Ehe jedoch in diesem erweiterten Zahlengebiet die Theorie der biquadratischen Reste selbst entwickelt werden kann, müssen in jenem die dieser Theorie vorangehenden Lehren der höhern Arithmetik, die bisher nur in Beziehung auf reelle Zahlen bearbeitet sind, an dieser Erweiterung Theil nehmen. Von diesen vorgängigen Untersuchungen können wir hier nur Einiges anführen. Der Verf. nennt jede Grösse $a+bi$, wo $a$ und $b$ reelle Grössen bedeuten, und $i$ der Kürze wegen anstatt $\surd-1$ geschrieben ist, eine complexe ganze Zahl, wenn zugleich $a$ und $b$ ganze Zahlen sind. Die complexen Grössen stehen also nicht den reellen entgegen, sondern enthalten diese als einen speciellen Fall, wo $b=0$, unter sich. Zur bequemern Handhabung war es erforderlich, mehrere auf die complexen Grössen sich beziehende Begriffsbildungen mit besondern Benennungen zu belegen, welche wir aber in dieser Anzeige zu umgehen suchen werden.

So wie in der Arithmetik der reellen Zahlen nur von zwei Einheiten, der positiven und negativen, die Rede ist, so haben wir in der Arithmetik der 
\SourceLeaf{182}%
complexen Zahlen vier Einheiten $+1$, $-1$, $+i$, $-i$. \emph{Zusammengesetzt} heisst eine complexe ganze Zahl, wenn sie das Product aus zwei von der Einheit verschiedenen ganzen Factoren ist; eine complexe Zahl hingegen, die eine \emph{solche} Zerlegung in Factoren nicht zulässt, heisst eine complexe Primzahl. So ist z. B. die reelle Zahl $3$, auch als complexe Zahl betrachtet, eine Primzahl, während $5$ als complexe Zahl zusammengesetzt ist $=(1+2i)(1-2i)$. Eben so wie in der höhern Arithmetik der reellen Zahlen spielen auch in dem erweiterten Felde dieser Wissenschaft die Primzahlen eine Hauptrolle.

Wird eine complexe ganze Zahl $a+bi$ als Modulus angenommen, so lassen sich $aa+bb$ unter sich nicht congruente, und nicht mehrere, complexe Zahlen aufstellen, von denen eine jede vorgegebene ganze complexe Zahl congruent sein muss, und die man ein vollständiges System incongruenter Reste nennen kann. Die sogenannten kleinsten und absolut kleinsten Reste in der Arithmetik der reellen Zahlen haben auch hier ihr vollkommenes Analogon. So besteht z. B. für den Modulus $1+2i$ das vollständige System der absolut kleinsten Reste aus den Zahlen $0$, $1$, $i$, $-1$ und $-i$. Fast die sämmtlichen Untersuchungen der vier ersten Abschnitte der \emph{Disquisitiones Arithmeticae} finden mit einigen Modificationen, auch in der erweiterten Arithmetik ihren Platz. Das berühmte \textsc{Fermat}sche Theorem z. B. nimmt hier folgende Gestalt an: Wenn $a+bi$ eine complexe Primzahl ist, und $k$ eine durch jene nicht theilbare complexe Zahl, so ist immer $k^{aa+bb-1}\equiv1$ für den Modulus $a+bi$. Ganz besonders merkwürdig ist es aber, dass das Fundamentaltheorem für die quadratischen Reste in der Arithmetik der complexen Zahlen sein vollkommenes, nur hier noch einfacheres, Gegenstück hat; sind nemlich $a+bi$, $A+Bi$ complexe Primzahlen, so dass $a$ und $A$ ungerade, $b$ und $B$ gerade sind, so ist die erste quadratischer Rest der zweiten, wenn die zweite quadratischer Rest der ersten ist, hingegen die erste quadratischer Nichtrest der zweiten, wenn die zweite quadratischer Nichtrest der ersten ist.

Indem die Abhandlung nach diesen Voruntersuchungen zu der Lehre von den biquadratischen Resten selbst übergeht, wird zuvörderst anstatt der blossen Unterscheidung zwischen biquadratischen Resten und Nichtresten eine Vertheilung der durch den Modulus nicht theilbaren Zahlen in vier Classen festgesetzt. Ist nemlich der Modulus eine complexe Primzahl $a+bi$, wo immer $a$ ungerade, $b$ gerade vorausgesetzt, und der Kürze wegen $p$ statt $aa+bb$ geschrieben wird, und $k$ eine complexe durch $a+bi$ nicht theilbare Zahl, so wird allemal $k^{\frac14(p-1)}$
\SourceLeaf{183}%
einer der Zahlen $+1,+i,-1,-i$ congruent sein, und dadurch eine Vertheilung sämmtlicher durch $a+bi$ nicht theilbarer Zahlen in vier Classen begründet, denen der Reihe nach der biquadratische Character $0,1,2,3$ beigelegt wird. Offenbar bezieht sich der Character $0$ auf die biquadratischen Reste, die übrigen auf die biquadratischen Nichtreste, und zwar so, dass dem Character $2$ zugleich quadratische\textperiodcentered\ Reste, den Charactern $1$ und $3$ hingegen quadratische Nichtreste entsprechen.

Man erkennt leicht, dass es hauptsächlich darauf ankommt, diesen Character blos für solche Werthe von $k$ bestimmen zu können, die selbst complexe Primzahlen sind, und hier führt sogleich die Induction zu höchst einfachen Resultaten.

Wird zuerst $k=1+i$ gesetzt, so zeigt sich, dass der Character dieser Zahl allemal $\equiv\frac18(-aa+2ab-3bb+1)(\mathrm{mod.}\ 4)$ wird, und ähnliche Ausdrücke finden sich für die Fälle $k=1-i$, $k=-1+i$, $k=-1-i$.

Ist hingegen $k=\alpha+\beta i$ eine solche Primzahl, wo $\alpha$ ungerade und $\beta$ gerade ist, so ergibt sich durch die Induction sehr leicht ein dem Fundamentaltheorem für die quadratischen Reste ganz analoges Reciprocitätsgesetz, welches am einfachsten auf folgende Art ausgedrückt werden kann:

Wenn sowohl $\alpha+\beta-1$ als $a+b-1$ durch $4$ theilbar sind (auf welchen Fall alle übrigen leicht zurückgeführt werden können), und der Character der Zahl $\alpha+\beta i$ in Beziehung auf den Modulus $a+bi$ durch $\lambda$, hingegen der Character von $a+bi$ in Beziehung auf den Modulus $\alpha+\beta i$ durch $l$ bezeichnet wird: so ist $\lambda=l$, wenn zugleich eine der Zahlen $\beta,b$ (oder beide) durch $4$ theilbar ist. hingegen $\lambda=l\pm2$, wenn keine der Zahlen $\beta,b$ durch $4$ theilbar ist.

Diese Theoreme enthalten im Grunde alles Wesentliche der Theorie der biquadratischen Reste in sich: so leicht es aber war, sie durch Induction zu entdecken, so schwer ist es, strenge Beweise für sie zu geben, besonders für das zweite, das Fundamentaltheorem der biquadratischen Reste. Wegen des grossen Umfanges, zu welchem schon die gegenwärtige Abhandlung angewachsen ist, sah sich der Verfasser genöthigt, die Darstellung des Beweises für das letztere Theorem, in dessen Besitz er seit $20$ Jahren ist, für eine künftige dritte Abhandlung zurückzulassen. Dagegen ist in vorliegender Abhandlung noch der vollständige Beweis für das erstere die Zahl $1+i$ betreffende Theorem (von welchem die 
\SourceLeaf{184}%
anderen für $1-i$, $-1+i$, $-1-i$ abhängig sind) mitgetheilt, welcher schon einigen Begriff von der Verwicklung des Gegenstandes geben kann.

Wir haben nun noch einige allgemeine Anmerkungen beizufügen. Die Versetzung der Lehre von den biquadratischen Resten in das Gebiet der complexen Zahlen könnte vielleicht manchem, der mit der Natur der imaginären Grössen weniger vertraut und in falschen Vorstellungen davon befangen ist, anstössig und unnatürlich scheinen, und die Meinung veranlassen, dass die Untersuchung dadurch gleichsam in die Luft gestellt sei, eine schwankende Haltung bekomme, und sich von der Anschaulichkeit ganz entferne. Nichts würde ungegründeter sein, als eine solche Meinung. Im Gegentheil ist die Arithmetik der complexen Zahlen der anschaulichsten Versinnlichung fähig, und wenngleich der Verf. in seiner diesmaligen Darstellung eine rein arithmetische Behandlung befolgt hat, so hat er doch auch für diese die Einsicht lebendiger machende und deshalb sehr zu empfehlende Versinnlichung die nöthigen Andeutungen gegeben. welche für selbstdenkende Leser zureichend sein werden. So wie die absoluten ganzen Zahlen durch eine in einer geraden Linie unter gleichen Entfernungen geordnete Reihe von Punkten dargestellt werden, in der der Anfangspunkt die Zahl $0$, der nächste die Zahl $1$ u. s. w. vertritt; und so wie dann zur Darstellung der negativen Zahlen nur eine unbegrenzte Verlängerung dieser Reihe auf der entgegengesetzten Seite des Anfangspunkts erforderlich ist: so bedarf es zur Darstellung der complexen ganzen Zahlen nur des Zusatzes, dass jene Reihe als in einer bestimmten unbegrenzten Ebene befindlich angesehen, und parallel mit ihr auf beiden Seiten eine unbeschränkte Anzahl ähnlicher Reihen in gleichen Abständen von einander angenommen werde, so dass wir anstatt einer Reihe von Punkten ein System von Punkten vor uns haben, die sich auf eine zweifache Art in Reihen von Reihen ordnen lassen, und zur Bildung einer Eintheilung der ganzen Ebene in lauter gleiche Quadrate dienen. Der nächste Punkt bei $0$ in der ersten Nebenreihe auf der einen Seite der Reihe, welche die reellen Zahlen repräsentirt, bezieht sich dann auf die Zahl $i$, so wie der nächste Punkt bei $0$ in der ersten Nebenreihe auf der andern Seite auf $-i$ u. s. f. Bei dieser Darstellung wird die Ausführung der arithmetischen Operationen in Beziehung auf die complexen Grössen, die Congruenz, die Bildung eines vollständigen Systems incongruenter Zahlen für einen gegebenen Modulus u. s. f. einer Versinnlichung fähig, die nichts zu wünschen übrig lässt.

\SourceLeaf{185}%
Von der andern Seite wird hiedurch die wahre Metaphysik der imaginären Grössen in ein neues helles Licht gestellt.

Unsere allgemeine Arithmetik, von deren Umfang die Geometrie der Alten so weit überflügelt wird, ist ganz die Schöpfung der neuern Zeit. Ursprünglich ausgehend von dem Begriff der absoluten ganzen Zahlen hat sie ihr Gebiet stufenweise erweitert; zu den ganzen Zahlen sind die gebrochenen, zu den rationalen die irrationalen, zu den positiven die negativen, zu den reellen die imaginären hinzugekommen. Dies Vorschreiten ist aber immer anfangs mit furchtsam zögerndem Schritt geschehen. Die ersten Algebraisten nannten noch die negativen Wurzeln der Gleichungen falsche Wurzeln, und sie sind es auch, wo die Aufgabe, auf welche sie sich beziehen, so eingekleidet vorgetragen ist, dass die Beschaffenheit der gesuchten Grösse kein Entgegengesetztes zulässt. Allein so wenig man in der \emph{Allgemeinen} Arithmetik Bedenken hat, die gebrochenen Zahlen mit aufzunehmen, obgleich es so viele zählbare Dinge gibt, wobei eine Bruchzahl ohne Sinn ist, eben so wenig durften in jener den negativen Zahlen gleiche Rechte mit den positiven deshalb versagt werden, weil unzählige Dinge kein Entgegengesetztes zulassen: die Realität der negativen Zahlen ist hinreichend gerechtfertigt, da sie in unzähligen andern Fällen ein adäquates Substrat finden. Darüber ist man nun freilich seit langer Zeit im Klaren: allein die den reellen Grössen gegenübergestellten imaginären --- ehemals. und hin und wieder noch jetzt, obwohl unschicklich, \emph{unmögliche} genannt --- sind noch immer weniger eingebürgert als nur geduldet, und erscheinen also mehr wie ein an sich inhaltleeres Zeichenspiel, dem man ein denkbares Substrat unbedingt abspricht, ohne doch den reichen Tribut, welchen dieses Zeichenspiel zuletzt in den Schatz der Verhältnisse der reellen Grössen steuert, verschmähen zu wollen.

Der Verf. hat diesen hochwichtigen Theil der Mathematik seit vielen Jahren aus einem verschiedenen Gesichtspunkt betrachtet, wobei den imaginären Grössen eben so gut ein Gegenstand untergelegt werden kann, wie den negativen: es hat aber bisher an einer Veranlassung gefehlt, dieselbe öffentlich bestimmt auszusprechen, wenn gleich aufmerksame Leser die Spuren davon in der $1799$ erschienenen Schrift über die Gleichungen, und in der Preisschrift über die Umbildung der Flächen leicht wiederfinden werden. In der gegenwärtigen Abhandlung sind die Grundzüge davon kurz angegeben; sie bestehen in Folgendem.

Positive und negative Zahlen können nur da eine Anwendung finden, wo
\SourceLeaf{186}%
das gezählte ein Entgegengesetztes hat, was mit ihm vereinigt gedacht der Vernichtung gleich zu stellen ist. Genau besehen findet diese Voraussetzung nur da Statt, wo nicht Substanzen (für sich denkbare Gegenstände) sondern Relationen zwischen je zweien Gegenständen das gezählte sind. Postulirt wird dabei, dass diese Gegenstände auf eine bestimmte Art in eine Reihe geordnet sind z. B. $A,B,C,D\dotsfour$, und dass die Relation des $A$ zu $B$ als der Relation des $B$ zu $C$ u. s. w. gleich betrachtet werden kann. Hier gehört nun zu dem Begriff der Entgegensetzung nichts weiter als der \emph{Umtausch} der Glieder der Relation, so dass wenn die Relation (oder der Uebergang) von $A$ zu $B$ als $+1$ gilt, die Relation von $B$ zu $A$ durch $-1$ dargestellt werden muss. Insofern also eine solche Reihe auf beiden Seiten unbegrenzt ist, repräsentirt jede reelle ganze Zahl die Relation eines beliebig als Anfang gewählten Gliedes zu einem bestimmten Gliede der Reihe.

Sind aber die Gegenstände von solcher Art, dass sie nicht in Eine, wenn gleich unbegrenzte, Reihe geordnet werden können, sondern sich nur in Reihen von Reihen ordnen lassen, oder was dasselbe ist, bilden sie eine Mannigfaltigkeit von zwei Dimensionen; verhält es sich dann mit den Relationen einer Reihe zu einer andern oder den Uebergängen aus einer in die andere auf eine ähnliche Weise wie vorhin mit den Uebergängen von einem Gliede einer Reihe zu einem andern Gliede derselben Reihe, so bedarf es offenbar zur Abmessung des Ueberganges von einem Gliede des Systems zu einem andern ausser den vorigen Einheiten $+1$ und $-1$ noch zweier andern unter sich auch entgegengesetzten $+i$ und $-i$. Offenbar muss aber dabei noch postulirt werden, dass die Einheit $i$ allemal den Uebergang von einem gegebenen Gliede einer Reihe zu einem \emph{bestimmten} Gliede der unmittelbar angrenzenden Reihe bezeichne. Auf diese Weise wird also das System auf eine doppelte Art in Reihen von Reihen geordnet werden können.

Der Mathematiker abstrahirt gänzlich von der Beschaffenheit der Gegenstände und dem Inhalt ihrer Relationen; er hat es blos mit der Abzählung und Vergleichung der Relationen unter sich zu thun: insofern ist er eben so, wie er den durch $+1$ und $-1$ bezeichneten Relationen, an sich betrachtet, Gleichartigkeit beilegt, solche auf alle vier Elemente $+1$, $-1$, $+i$ und $-i$ zu erstrecken befugt.

Zur Anschauung lassen sich diese Verhältnisse nur durch eine Darstellung
\SourceLeaf{187}%
im Raume bringen, und der einfachste Fall ist, wo kein Grund vorhanden ist, die Symbole der Gegenstände anders als quadratisch anzuordnen, indem man nemlich eine unbegrenzte Ebene durch zwei Systeme von Parallellinien, die einander rechtwinklig durchkreuzen, in Quadrate vertheilt, und die Durchschnittspunkte zu den Symbolen wählt. Jeder solche Punkt $A$ hat hier vier Nachbaren, und wenn man die Relation des $A$ zu einem benachbarten Punkte durch $+1$ bezeichnet, so ist die durch $-1$ zu bezeichnende von selbst bestimmt, während man, welche der beiden andern man will, für $+i$ wählen, oder den sich auf $+i$ beziehenden Punkt nach Gefallen \emph{rechts} oder \emph{links} nehmen kann. Dieser Unterschied zwischen rechts und links ist, so bald man vorwärts und rückwärts \emph{in} der Ebene, und oben und unten in Beziehung auf die beiden Seiten der Ebene einmal (nach Gefallen) festgesetzt hat, \emph{in sich} völlig bestimmt, wenn wir gleich unsere Anschauung dieses Unterschiedes andern \emph{nur} durch Nachweisung an wirklich vorhandenen materiellen Dingen mittheilen können\begingroup\renewcommand{\thefootnote}{*)}\footnote{Beide Bemerkungen hat schon \textsc{Kant} gemacht, aber man begreift nicht, wie dieser scharfsinnige Philosoph in der ersteren einen Beweis für seine Meinung, dass der Raum \emph{nur} Form unserer äussern Anschauung sei, zu finden glauben konnte, da die zweite so klar das Gegentheil, und dass der Raum unabhängig von unserer Anschauungsart eine reelle Bedeutung haben muss, beweiset.}\endgroup. Wenn man aber auch über letzteres sich entschlossen hat, sieht man, dass es doch von unserer Willkür abhing, welche von den beiden in Einem Punkte sich durchkreuzenden Reihen wir als Hauptreihe, und welche Richtung in ihr man als auf positive Zahlen sich beziehend ansehen wollten; man sieht ferner, dass wenn man die vorher als $+i$ behandelte Relation für $+1$ nehmen will, man nothwendig die vorher durch $-1$ bezeichnete Relation für $+i$ nehmen muss. Das heisst aber, in der Sprache der Mathematiker, $+i$ ist mittlere Proportionalgrösse zwischen $+1$ und $-1$ oder entspricht dem Zeichen $\surd-1$: wir sagen absichtlich nicht \emph{die} mittlere Proportionalgrösse, denn $-i$ hat offenbar gleichen Anspruch. Hier ist also die Nachweisbarkeit einer anschaulichen Bedeutung von $\surd-1$ vollkommen gerechtfertigt, und mehr bedarf es nicht, um diese Grösse in das Gebiet der Gegenstände der Arithmetik zuzulassen.

Wir haben geglaubt, den Freunden der Mathematik durch diese kurze Darstellung der Hauptmomente einer neuen Theorie der sogenannten imaginären Grössen einen Dienst zu erweisen. Hat man diesen Gegenstand bisher aus einem falschen Gesichtspunkt betrachtet und eine geheimnissvolle Dunkelheit dabei 
\SourceLeaf{188}%
gefunden, so ist dies grossentheils den wenig schicklichen Benennungen zuzuschreiben. Hätte man $+1$, $-1$, $\surd-1$ nicht positive, negative, imaginäre (oder gar unmögliche) Einheit, sondern etwa directe, inverse, laterale Einheit genannt, so hätte von einer solchen Dunkelheit kaum die Rede sein können. Der Verf. hat sich vorbehalten, den Gegenstand, welcher in der vorliegenden Abhandlung eigentlich nur gelegentlich berührt ist, künftig vollständiger zu bearbeiten, wo dann auch die Frage, warum die Relationen zwischen Dingen, die eine Mannigfaltigkeit von mehr als zwei Dimensionen darbieten, nicht noch andere in der allgemeinen Arithmetik zulässige Arten von Grössen liefern können, ihre Beantwortung finden wird.
\BookRule{.28\linewidth}
